%\documentclass[twocolumn]{article}\
\documentclass{article}
\usepackage{amsmath, amssymb, cancel, mathtools, bm}
\usepackage[left=.5in, right=.5in, top=1in, bottom=1in]{geometry}
\usepackage[most]{tcolorbox}
\usepackage[skip=1em,indent=0pt]{parskip}
\setlength{\parindent}{0pt}
\newcommand{\statvec}[1]{\underset{\sim}{\bm{#1}}} % Vector symbol (tilde under vector)
\DeclareMathOperator{\EX}{\mathbb{E}} % Expected Value symbol
\DeclareMathOperator{\ext}{\EX_{\theta}} % Expected Value symbol
\DeclareMathOperator{\vart}{Var_{\theta}} % Expected Value symbol
\newcommand{\indep}{\perp\!\!\!\!\perp} % Independence symbol
\newcommand{\real}{\mathbb{R}} % Simplified 'Reals' indicator
\newcommand{\pto}{\overset{P}{\to}}
\newcommand{\asto}{\overset{a.s.}{\to}}
\newcommand{\rthto}{\overset{L^r}{\to}}
\newcommand{\dto}{\overset{D}{\to}}
% Force all aggregate symbols to always put values above/below
\let\Oldint=\int
\let\Oldsum=\sum
\let\Oldprod=\prod
\let\Oldbigcup=\bigcup
\let\Oldbigcap=\bigcap
\let\Oldlim=\lim
\renewcommand{\int}{\Oldint\limits} 
\renewcommand{\sum}{\Oldsum\limits} 
\renewcommand{\prod}{\Oldprod\limits} 
\renewcommand{\bigcup}{\Oldbigcup\limits} 
\renewcommand{\bigcap}{\Oldbigcap\limits} 
\renewcommand{\lim}{\Oldlim\limits} 
\newcommand{\infint}{\int_{-\infty}^{\infty}}
\newcommand{\iiddist}{\overset{\mathrm{iid}}{\sim}}
\newcommand{\norm}[1]{\left\lVert#1\right\rVert}
\newcommand{\boldred}[1]{\textbf{\textcolor{red}{#1}}}



\begin{document}

%%%%%%%%%%%%%%%%%%%%%%%%%%%%%%%%%%%%%%%%%%%%%%%%%%%%%%%%%%%%%%%%%%%%%%%%%%%

\underline{Throughout everything, justify my steps!}

\section{Method of Moments}

Dth Population Mean: $E(X^d) = \mu_d(\theta)$

Dth Sample Moment: $\bar{\mu}_d = \frac{1}{n}\sum x_i^d$

Set them equal to eachother, and solve the system of equations: $\begin{matrix} \bar{\mu}_1 = \mu_1(\theta) \\ \bar{\mu}_2 = \mu_2(\theta) \\ \vdots \end{matrix}$

\section{MLE}

Find the derivative, set equal to 0. \underline{Include the 2nd derivative!} (make sure it's negative)

Hessian: Determinants should alternate negative and positive

$\hat{\theta}$ is MLE of $\theta$: $\tau(\hat{\theta})$ is MLE for $\tau(\theta)$

\section{Loss}: $L(\theta, T)$

Squared Error: $(T-\theta)^2$  - Bayes Rule: Posterior Mean

Absolute Error: $|T - \theta|$ - Bayes Rule: Posterior Median

Zero-One: $\begin{cases} 0 & T = \theta \\ 1 & T \neq \theta \end{cases}$ - Bayes Rule: Posterior Mode

Risk: $\EX[L(\theta, T)]$

Bayes Rule: $\EX[L(\theta, T)|x] = \int_{\mathbb{X}} L(\theta, T)\pi(\theta|x)d\theta$

Goal in the process: get the Error into one of the above forms, then absorb the rest into the posterior in the integral.

"Better" Risk: $R_{T_1}(\theta) \le R_{T_2}(\theta) \ \forall \theta, R_{T_1}(\theta) < R_{T_2}(\theta) \ \exists \theta$

\section{CRLB}

$CRLB_{\theta} = \frac{(\tau'(\theta))^2}{I_n(\theta)}$. Conditions: Bounds for $f(x|\theta)$ not dependent on $\theta$. The Fisher information needs to exist and be finite.

If the CRLB is reached by an estimator, it is the MVUE.

$I_n(\theta) = -\EX[\frac{d^2}{d\theta^2}ln(f(x|\theta))]$

Sam's Corrolary: Iff $\frac{\partial}{\partial \theta} ln(L(\theta|x)) = a(\theta)[W(x) - \tau(\theta)]$, then $W(x)$, estimating $\tau(\theta)$, attains the CRLB (and thus is the MVUE).

\section{Sufficiency}

$(x_1, ..., x_n|S)$ does not depend on $\theta$ (ie. $\frac{P(x_1, ..., x_n)}{P(S)} = ...$ which does not depend on $\theta$)

Factorization: $f(x|\theta) = g(S, \theta)h(x)$. Any 1:1 function of S is also sufficient.

Rao-Blackwell: $T$ is unbiased for $\tau(\theta)$, $S$ is sufficient: $\tau^* = \EX[T|S], Var(T^*) \le Var(T)$

Ancillary: Statistic whose distribution does not depend on $\theta$

\section{Completeness}

$T$ is complete if, for $g(T)$, $\EX[g(T)] = 0 \ \forall \theta \in \Theta \implies P(g(T) = 0) = 1 \ \forall \theta \in \Theta$

Basically, $g(T)$ is the only part that could be 0

Exponential Family: $f(x|\theta) = h(x)c(\theta)e^{\sum+{j=1}^k w_j(\theta)t_j(x)}: \ T(X) = (t_1(x), ..., t_2(x))$ is complete as long as $\Theta$ contains an open set in $\real^k$.

Basu's: S is complete and sufficient, T is ancillary, S and T are independent

\section{Lehman-Scheffe}

S is complete and sufficient. If $T = g(S)$ is unbiased for $\tau(\theta), T$ is the MVUE 

Steps: 

1. Find complete and sufficient Statistic

2. Check bias, adjust Statistic

\section{Other}

$\frac{d}{dx} \int_{g(x)}^{h(x)} f(t)dt = f(h(x))h'(x) - f(g(x))g'(x)$

$MSE = Var + Bias^2$

\end{document}